\documentclass[12pt,a4paper]{article}
\usepackage[T1]{fontenc}
\usepackage[utf8]{inputenc}
\usepackage[brazil]{babel}

\title{Mundo dos Blocos de Tamanho Variável via SAT Solver}
\author{Equipe}
\date{Outubro de 2026}

\begin{document}
\maketitle

\section{Introdução ao problema}
% TODO: apresentar o objetivo e os três cenários.

\section{Descrição formal do mundo dos blocos de tamanho variável}
% TODO: LPO, predicados, ações, efeitos, estabilidade e ordem parcial.

\section{Codificação CNF}
% TODO: variáveis proposicionais e grupos de cláusulas.

\section{Exemplos dos três cenários codificados em lógica proposicional}
% TODO: estados iniciais, metas e exemplos de cláusulas.

\section{Mapeamento da descrição formal para o código}
% TODO: relacionar cada regra ao gerador Python.

\section{Execução passo a passo}
% TODO: geração da CNF, mapa e execução do SAT solver.

\section{Interpretação da saída do SAT solver}
% TODO: decodificação dos resultados e comparação com os planos manuais.

\end{document}
